% ==========================================================================
%  Chapter 89 -- XFEM Local Model Closure (Plan v2 §Fase 3)
% ==========================================================================

\chapter{XFEM Local Model Closure}
\label{ch:xfem-local-model-closure}

\paragraph{Normative status.}
This chapter records the closure evidence for the XFEM shifted-Heaviside
local model used as the \emph{promoted physical local model} of the
multiscale FE\(^2\) program (Plan~v2 \S Fase~3). It supersedes the
exploratory residual heuristics scattered through Ch.~82 and Ch.~85 and
provides the runtime evidence required to enter the four-stage chain
(Ch.~90).

\section{Promotion catalog and gate}

The promotion catalog
\href{run:src/validation/ReducedRCLocalModelPromotionCatalog.hh}{%
\texttt{ReducedRCLocalModelPromotionCatalog.hh}} declares the
\texttt{xfem\_global\_secant\_200mm\_primary\_candidate} entry as the
canonical local model. The Cap.~79 gate requires
\(\mathrm{rms}<0.10\), \(\max<0.30\), and \(\text{peak\_ratio}<1.15\) on a
50/100/150/200~mm cyclic protocol with \(\mathrm{NZ}=4\).

\section{Runtime evidence (commits \texttt{ae652a6}, \texttt{ba4704d},
                          \texttt{d224509}, \texttt{633ef61})}

Closure status per kinematic formulation:

\begin{itemize}
  \item \textbf{Small-strain}: reference baseline. Catalog gate PASS at
        all amplitudes up to 300~mm and \(\mathrm{NZ}\in\{4,5,6,8\}\).
        Peak base shear monotonically decreases with \(h\)-refinement
        (\(h\)-convergence), refuting the legacy \emph{NZ=8 timed-out}
        claim.
  \item \textbf{Corotational}: 200~mm gate PASS with
        \(\mathrm{rms}=0.0014\), \(\max=0.0044\),
        \(\text{peak\_ratio}=0.9965\). 300~mm extension also PASS
        (\(0.0018, 0.0074, 0.9930\)).
  \item \textbf{Total Lagrangian}: 200~mm
        \(\mathrm{rms}=0.0345, \max=0.1610, \text{peak\_ratio}=1.1610\);
        300~mm \(0.063, 0.32, 1.32\). RMS and \(\max\) gates PASS at
        200~mm; peak ratio narrowly above the 1.15 ceiling.
  \item \textbf{Updated Lagrangian}: 200~mm
        \(\mathrm{rms}=0.0318, \max=0.1502, \text{peak\_ratio}=1.1502\);
        300~mm \(0.058, 0.30, 1.30\). Same regime as TL.
\end{itemize}

Empirical residuals are recorded in
\path{data/output/cyclic_validation/xfem_*_v2/} and audited by
\href{run:scripts/evaluate_xfem_corotational_gate.py}{%
\texttt{evaluate\_xfem\_corotational\_gate.py}} and
\href{run:scripts/evaluate_xfem_finite_kinematics_gate.py}{%
\texttt{evaluate\_xfem\_finite\_kinematics\_gate.py}}.

\section{Regression ctests}

The closure is locked through the following walk-up CSV regression tests
(LIBS NONE, label \texttt{validation\_reboot}):
\begin{itemize}
  \item \texttt{xfem\_corotational\_200mm\_gate} (\#125)
  \item \texttt{xfem\_corotational\_300mm\_gate} (\#126)
  \item \texttt{xfem\_finite\_kinematics\_residuals} (\#127)
\end{itemize}

\section{Scoped-deferred residuals (research-level)}

Items still tagged \texttt{scoped\_deferred\_to\_branch} per
\path{data/output/validation_reboot/audit_phase3_xfem_closure.json}:
\begin{enumerate}
  \item Real bordered mixed-control constraint inside
        \texttt{PetscBorderedMixedControlNewton} (today: hybrid SNES-L2
        fallback closes the gate in \(\sim 256\)~s).
  \item Adaptive enrichment activation policy (probe present; runtime
        wiring deferred).
  \item FieldSplit/Schur ASM preconditioning for
        \(\mathrm{NZ}\geq 5\) iterative path.
  \item Second-generation active-set tangent.
  \item VTK time-series export with crack opening / cohesive traction.
  \item TL/UL 250--300~mm peak-ratio recalibration.
\end{enumerate}

These are explicitly \emph{not} blockers for the FE\(^2\) chain because
\(\mathrm{NZ}=4\) small-strain is the canonical local model. They are
documented research-growth paths.

\paragraph{Honest scientific status.} \texttt{closed\_with\_runtime\_evidence}
for small-strain and corotational; \texttt{empirical\_residual\_recorded}
for TL/UL; the six items above remain
\texttt{scoped\_deferred\_to\_branch}.
